\section{Discussion}
\label{sec:discussion}

\subsection{The Coupling Disparity}
\label{sec:disc:disparity}

% Deliberately crude stick-figure diagram comparing excitation pathways
% Intended to be charming rather than professional
\begin{figure}[htbp]
\centering
\begin{tikzpicture}[
    thick,
    every node/.style={font=\footnotesize},
    wave/.style={decorate, decoration={snake, amplitude=2pt, segment length=6pt}},
]

% =========================================================
% LEFT: Airborne pathway (no effect)
% =========================================================
\begin{scope}[shift={(-3.5,0)}]
  % Speaker
  \draw[fill=gray!30] (-2.8,1.5) rectangle (-2.3,0.5);
  \draw (-2.3,0.5) -- (-1.8,0) -- (-1.8,2) -- (-2.3,1.5);
  \node[above] at (-2.3,2.1) {\textbf{120\,dB}};

  % Sound waves
  \draw[wave, ->] (-1.7,1.0) -- (-0.5,1.0);

  % Stick figure - standing, unbothered
  % Head
  \draw (0,2.2) circle (0.25);
  % Smiley face
  \draw (-0.08,2.28) arc (220:320:0.1);
  \fill (-0.08,2.28) circle (0.03);
  \fill (0.08,2.28) circle (0.03);
  % Body
  \draw (0,1.95) -- (0,1.0);
  % Abdomen (circle to represent cavity)
  \draw[dashed] (0,1.3) ellipse (0.25 and 0.2);
  \node[right, font=\tiny, fill=white, inner sep=1pt] at (0.45,1.45)
    {$\xi = \SI{0.014}{\micro\metre}$};
  % Arms
  \draw (0,1.7) -- (-0.4,1.3);
  \draw (0,1.7) -- (0.4,1.3);
  % Legs
  \draw (0,1.0) -- (-0.3,0.3);
  \draw (0,1.0) -- (0.3,0.3);
  % Ground
  \draw[thin] (-1.0,0.3) -- (1.0,0.3);

  % Thought bubble
  \node[draw, rounded corners, fill=white, font=\tiny, align=center]
    at (0.95,2.85) {perfectly\\fine};

  \node[below, font=\small\bfseries, align=center] at (0,-0.02) {(a) Airborne\\[1pt]``Nothing happens''};
\end{scope}

% =========================================================
% RIGHT: Mechanical pathway (considerable effect)
% =========================================================
\begin{scope}[shift={(3.5,0)}]
  % Vibrating platform
  \draw[fill=gray!30] (-1.0,0.3) rectangle (1.0,0.5);
  \draw[wave] (-1.0,0.1) -- (1.0,0.1);
  \node[below, font=\tiny, fill=white, inner sep=1pt] at (0,0.03)
    {$a = \SI{0.5}{\metre\per\second\squared}$};

  % Up arrows (vibration)
  \draw[->, thick] (-0.5,0.5) -- (-0.5,0.8);
  \draw[->, thick] (0.5,0.5) -- (0.5,0.8);
  \draw[->, thick] (0,0.5) -- (0,0.8);

  % Stick figure - distressed
  % Head
  \draw (0,2.2) circle (0.25);
  % Distressed face
  \draw (-0.08,2.28) arc (140:40:0.1);
  \fill (-0.08,2.28) circle (0.03);
  \fill (0.08,2.28) circle (0.03);
  \draw (0,2.12) arc (220:320:0.06);
  % Body (slightly wavy to suggest vibration)
  \draw (0,1.95) -- (0.05,1.5) -- (-0.05,1.2) -- (0,0.8);
  % Abdomen (larger, distressed)
  \draw[dashed, thick] (0,1.3) ellipse (0.3 and 0.25);
  \node[right, font=\tiny, fill=white, inner sep=1pt] at (0.48,1.48)
    {$\xi = \SI{0.9}{\milli\metre}$};
  % Arms (thrown up in distress)
  \draw (0,1.7) -- (-0.5,2.2);
  \draw (0,1.7) -- (0.5,2.2);
  % Legs
  \draw (0,0.8) -- (-0.3,0.5);
  \draw (0,0.8) -- (0.3,0.5);

  % Distress indicators
  \node[font=\tiny, rotate=15] at (-0.7,2.5) {!};
  \node[font=\tiny, rotate=-10] at (0.7,2.5) {!};

  % The brown note consequence (subtle)
  \draw[->, decorate, decoration={snake, amplitude=1pt, segment length=3pt}]
    (0.15,0.7) -- (0.4,0.4);
  \node[right, font=\tiny, align=left, fill=white, inner sep=1pt] at (0.58,0.55)
    {bowel\\[-1pt]response};

  \node[below, font=\small\bfseries, align=center] at (0,-0.02)
    {(b) Whole-body vibration\\[1pt]``The resonance is real''};
\end{scope}

% =========================================================
% Centre annotation
% =========================================================
\draw[<->, very thick] (-0.9,2.75) -- (0.9,2.75);
\node[above, font=\small, align=center, fill=white, inner sep=1pt] at (0,2.8)
  {$\times\, 66{,}000$};

\end{tikzpicture}
\caption{Schematic comparison of the two excitation pathways.  (a)~Airborne
infrasound at \SI{120}{\decibel} SPL produces wall displacements of
\SI{0.014}{\micro\metre}---far below any conceivable physiological threshold.
The subject is entirely unperturbed.  (b)~Whole-body vibration at occupational
levels produces displacements roughly $6.6 \times 10^4$ times larger, comfortably
exceeding mechanotransduction thresholds and---as the occupational health
literature attests---occasionally exceeding the subject's composure as well.
Diagram intentionally not to scale, anatomically, or artistically.}
\label{fig:stick_figure}
\end{figure}


Figure~\ref{fig:stick_figure} summarises this contrast schematically.  The
central result of this study is the separation between the airborne and
mechanical transfer functions, not any single exposure-pair ratio.
Canonically, $H_\mathrm{mech}/H_\mathrm{air}
\approx 6.5 \times 10^6\,\si{\pascal\second\squared\per\metre}$ before the
modal-participation correction, and Monte Carlo analysis shows that this
transfer-function contrast remains decisive across physiologically plausible
parameter ranges.  For the commonly quoted reference pair
(\SI{120}{\decibel} SPL airborne excitation and
$a_\mathrm{rms} = \SI{0.1}{\metre\per\second\squared}$ RMS whole-body
vibration), the same transfer-function separation implies the SDOF upper-bound
displacement ratio $\mathcal{R}_\mathrm{full} \approx 3.3 \times 10^4$;
including the canonical participation factor $\Gamma_2 \approx 0.48$ reduces
this to $\sim 1.6 \times 10^4$
(Tables~\ref{tab:airborne}--\ref{tab:mechanical}).  The physical origin of
this disparity lies in two independent mechanisms: a geometric filtering
penalty that suppresses modal-pressure content, and an impedance mismatch
that limits energy transfer.

First, the air--tissue impedance mismatch is severe.  The specific acoustic
impedance of air is $Z_{\mathrm{air}} = \rho_{\mathrm{air}} c_{\mathrm{air}}
\approx \SI{420}{\pascal\second\per\metre}$, while that of soft tissue is
$Z_{\mathrm{tissue}} \approx \SI{1.6e6}{\pascal\second\per\metre}$, giving
an intensity transmission coefficient
$T_I = 4 Z_{\mathrm{air}} Z_{\mathrm{tissue}} /
(Z_{\mathrm{air}} + Z_{\mathrm{tissue}})^2 \approx 1.1 \times 10^{-3}$.
At the body surface, approximately \SI{99.9}{\percent} of the incident
acoustic power is reflected, before the still smaller mode-specific resonant
absorption is accounted for.  This general suppression of transmission across large biological
impedance contrasts is consistent with modern frequency-dependent analyses of
bone--soft tissue interfaces \cite{Gupta2021}.

Second, even the transmitted pressure must couple to the spatial pattern of
a flexural mode.  In the long-wavelength (Rayleigh) limit, $ka \ll 1$, the
$n$-th order component of the incident plane wave scales as $(ka)^n$
\cite{Junger1986}.  At the $n{=}2$ resonance
($f_2 \approx \SI{4}{\hertz}$) with
$R_{\mathrm{eq}} = \SI{0.157}{\metre}$, we obtain
$ka = 2\pi f R / c_{\mathrm{air}} \approx 0.0114$ and therefore
$(ka)^2 \approx 1.3 \times 10^{-4}$ for the $n = 2$ flexural mode
(equation~\ref{eq:peff2}).  Physically, the incident wavelength
($\lambda \approx \SI{86}{\metre}$) is nearly 270 times the body diameter;
the pressure field is almost perfectly uniform over the body surface, and
the quadrupolar ($n = 2$) pressure gradient that could drive an
oblate--prolate deformation is vanishingly small.

The energy budget analysis (\S\ref{sec:coupling}) confirms this picture from
an independent perspective.  The radiation damping ratio in air is
$\zeta_{\mathrm{rad}} \approx 10^{-15}$, whereas the structural damping
ratio is $\zeta_{\mathrm{struct}} = 0.125$, so the absorption efficiency
$\zeta_{\mathrm{rad}} / \zeta_{\mathrm{total}} \sim 10^{-14}$, giving a
Breit--Wigner absorption efficiency
$4\zeta_\mathrm{rad}\zeta_\mathrm{struct}/\zeta_\mathrm{total}^2 \sim 3 \times 10^{-14}$.  The shell
is, for practical purposes, acoustically transparent at infrasonic
frequencies.  By contrast, the upper-bound mechanical model assumes rigid
skeletal transmission: the bone--peritoneum impedance contrast is orders of
magnitude smaller than the air--tissue mismatch, and the base-excitation
drives the $n = 2$ mode via equation~\ref{eq:Hrel} with a modal
participation factor $\Gamma_2 \approx 0.48$
(\S\ref{sec:coupling_ratio}).

This disparity resolves an apparent paradox in the literature.  Whole-body
vibration at \SIrange{4}{8}{\hertz} is well documented to cause
gastrointestinal distress---including nausea, abdominal pain, and in
extreme cases, incontinence \cite{Griffin1990,mansfield2005human,Seidel1986,Ishitake2002,Krajnak2018}---whereas
airborne infrasound studies at levels up to \SI{150}{\decibel} report no
comparable effects \cite{Moller2004,Leventhall2007}.  The present modal
analysis shows that both observations are consistent with the \emph{same}
underlying resonance.  The two excitation pathways differ not in the
resonance they address, but in the efficiency with which they address it.

\subsection{Implications for the Brown Note Hypothesis}
\label{sec:disc:brownnote}

The results permit a quantitative evaluation of the ``brown note''
hypothesis.  The hypothesis contains one correct element: abdominal resonance in the
\SIrange{4}{10}{\hertz} band is physically real, well-supported by the
parametric study (\S\ref{sec:results}), and consistent with ISO~2631
transmissibility data \cite{iso2631,Kitazaki1998}.  Independent
finite-element studies of the seated human body now recover dominant internal
and abdominal vibration modes in the \SIrange{4}{8}{\hertz} range
\cite{Dong2024}.  The flexural
$n = 2$ mode at \SI{4.0}{\hertz} (for $E = \SI{0.1}{\mega\pascal}$, relaxed
musculature) falls squarely within the frequency band where occupational
exposure studies report gastrointestinal symptoms.

Where the hypothesis fails is in the \emph{excitation pathway}.  At
\SI{120}{\decibel} SPL---a level exceeding the conventional threshold of
discomfort and well above any
environmental or entertainment exposure---the predicted airborne wall
displacement is the energy-consistent value
\SI{0.028}{\micro\metre} (Table~\ref{tab:airborne}), roughly
two orders of magnitude below the \SIrange{0.5}{2.0}{\micro\metre} cellular
activation range of mechanosensitive PIEZO1 and PIEZO2 ion channels
\cite{coste2010piezo1,Lewis2015,Xiao2024} (a benchmark derived from
cell-culture experiments; see Limitation~9).  In the modern mechanotransduction
literature, the more relevant gastrointestinal transducer is probably PIEZO2:
somatosensory-neuron PIEZO2 signalling has now been shown to regulate gut
transit directly \cite{ServinVences2023}.  Reaching the \SI{1}{\micro\metre}
cellular benchmark on the energy-consistent basis requires approximately
\SI{151}{\decibel} (Table~\ref{tab:uq_summary} gives a median of
\SI{152}{\decibel} with a 90\% credible interval of
\SIrange{145}{156}{\decibel}),
a level associated with severe acoustic hazards and encountered only in
explosive blasts and rocket launches.  Even at
this extreme, the displacement remains three
orders of magnitude below what routine whole-body vibration
($a_{\mathrm{rms}} = \SI{0.5}{\metre\per\second\squared}$) produces
(Table~\ref{tab:mechanical}).

The model also provides the first quantitative explanation for \emph{why}
whole-body vibration at \SIrange{4}{8}{\hertz} is so effective at producing
gastrointestinal effects: it excites a flexural resonance of the
fluid-filled abdominal cavity, producing millimetre-scale tissue
displacements that far exceed cellular mechanotransduction benchmarks.  The
epidemiological pattern documented by ISO~2631---peak symptom incidence at
\SIrange{4}{8}{\hertz}, with a strong dependence on acceleration
amplitude---is precisely what a resonance-driven mechanism predicts.

\subsection{Alternative Coupling Pathways}
\label{sec:disc:alternatives}

The whole-cavity resonance model of \S\ref{sec:formulation} treats the
abdominal wall as a homogeneous fluid-filled shell.  In reality, alternative
local coupling mechanisms may also exist, notably compressible intestinal gas
pockets and intermittent orifice coupling through the upper
gastrointestinal tract.  These mechanisms would operate through local
compliance and transient boundary conditions rather than through the global
$n=2$ cavity resonance analysed here.  We therefore treat them only as
hypotheses for future study, not as co-equal conclusions of the present paper:
the result established here is the severe weakness of airborne coupling to the
global abdominal flexural mode relative to mechanical excitation.  Assessing
whether any local pathway can become biologically relevant under extreme
conditions would require a dedicated model that resolves gas--tissue
interactions, anatomical intermittency, and local mechanotransduction.

\subsection{Clinical and Occupational Relevance}
\label{sec:disc:clinical}

The transfer-function disparity has direct implications for occupational exposure
standards.  ISO~2631-1 specifies frequency-weighted acceleration limits for
whole-body vibration, with the highest sensitivity weighting in the
\SIrange{4}{8}{\hertz} band---precisely the range where our model predicts
the $n = 2$ flexural resonance \cite{iso2631}, and consistent with
von~Gierke and Brammer's compilation of organ resonance
data~\cite{vonGierke2002}.  More recent occupational reviews continue to link
whole-body vibration exposure with adverse health outcomes across multiple
organ systems \cite{Krajnak2018}.  The EU Physical Agents
(Vibration) Directive~2002/44/EC codifies two thresholds: a daily exposure
action value of \SI{0.5}{\metre\per\second\squared} $A(8)$ and an exposure
limit value of \SI{1.15}{\metre\per\second\squared} $A(8)$.  (The $A(8)$
action and limit values are frequency-weighted eight-hour equivalent
accelerations; we use them here as illustrative equivalent harmonic
amplitudes at the $n{=}2$ resonance frequency, not as direct predictions
under the Directive's time-weighted assessment protocol.)  At the action
value, our model predicts a relative abdominal wall displacement of
\SI{4.6}{\milli\metre} at resonance (Table~\ref{tab:mechanical})---more than
$4{,}000\times$ the PIEZO cellular activation benchmark.  The present analysis
therefore provides a \emph{mechanistic rationale} for why these particular
frequency bands and amplitude thresholds were empirically selected: they
correspond to conditions under which flexural resonance of the abdominal
cavity produces tissue strains large enough to activate
mechanotransduction pathways.

For airborne infrasound, ISO~7196 defines G-weighting but prescribes no
dose--response limits.  Our analysis supports this regulatory asymmetry:
at environmental and occupational SPL ($\leq \SI{120}{\decibel}$), airborne
coupling to abdominal flexural modes is negligible, and exposure limits
based on abdominal resonance effects are unnecessary at these levels.

A potentially valuable extension of this framework is toward personalised
vibration exposure assessment.  The parametric study (\S\ref{sec:results})
shows that the $n = 2$ eigenfrequency depends sensitively on body habitus:
larger abdominal radius (higher BMI) reduces the resonance frequency, while
greater muscular tone (higher effective $E$) increases it.  In principle,
imaging-derived geometric and viscoelastic parameters \cite{Sack2023} combined with
the analytical expressions of \S\ref{sec:formulation} could yield
subject-specific resonance frequency estimates, enabling targeted vibration
exposure limits for workers whose body habitus places their resonance
frequency in a particularly hazardous band.

\subsection{Broader Applications of the Analytical Framework}
\label{sec:disc:broader}

The brown note provided the motivation for this work, but the analytical
framework---a fluid-filled viscoelastic shell, a modal decomposition into
breathing and flexural families, and a systematic comparison of airborne
versus mechanical coupling efficiency---is considerably more general.  Any
biological or engineered structure that can be idealised as a compliant
enclosure surrounding a dense fluid falls within the scope of the model
developed in \S\ref{sec:formulation}.  The two key dimensionless groups that
govern the physics, the Helmholtz number $ka$ and the impedance ratio
$Z_{\mathrm{wall}}/Z_{\mathrm{fluid}}$, transfer directly; what changes
between applications is the frequency range, the material properties, and
the geometric aspect ratio of the shell.  Potential applications include
primary blast injury modelling---where $ka$ approaches unity and the
geometric filtering penalty largely disappears
\cite{Stuhmiller2008,Owers2011}---swim-bladder target-strength prediction
in fisheries acoustics \cite{Feuillade1998,Love1978}, and therapeutic
ultrasound dosimetry where coupling must be maximised rather than shown to be
negligible \cite{terHaar2007}.  These extensions are beyond the scope of the
present work but follow directly from the same modal decomposition and
coupling-comparison methodology.

\subsection{Nonlinear Amplitude Effects}
\label{sec:disc:nonlinear}

The linear model of \S\ref{sec:formulation} is valid only for small
displacements relative to the wall thickness.  At the large amplitudes
associated with whole-body vibration near the EU exposure limit, geometric
nonlinearity from membrane stretching and fluid volume conservation may
stiffen the response, further reducing actual displacement below the linear
prediction.  Such hardening behaviour is characteristic of thin-shell
vibration at moderate amplitudes \cite{Alijani2014} and would shift the
resonance peak to higher frequency while reducing its amplitude---effects that
are conservative for the brown note argument, since they widen the gap between
prediction and any putative airborne effect.  The airborne displacement
($\sim\SI{0.028}{\micro\metre}$) remains deep in the linear regime regardless.
A detailed nonlinear analysis, including the amplitude-dependent frequency
shift and hysteretic jump phenomena expected at occupational vibration
levels, is deferred to future work.

\subsection{Limitations}
\label{sec:disc:limitations}

We identify the following limitations, grouped by their likely effect on
the conclusions.

\begin{enumerate}

  \item \textbf{Linear theory.}  The Love--Kirchhoff shell formulation
    assumes small displacements.  As discussed in
    \S\ref{sec:disc:nonlinear}, occupational vibration amplitudes likely
    exceed the linear regime, and the resulting hardening nonlinearity would
    reduce peak displacement below the linear prediction.  Because the airborne
    displacement remains deep in the linear regime, nonlinearity
    \emph{reduces} the coupling ratio somewhat but does not alter the
    qualitative conclusion.

  \item \textbf{Homogeneous, isotropic wall.}  The abdominal wall is a
    layered anisotropic composite comprising peritoneum, transversalis
    fascia, internal oblique muscle, external oblique muscle, subcutaneous
    fat, and skin.  The homogenisation error in effective bending stiffness
    is expected to be modest compared with inter-subject variability in
    wall thickness and elastic modulus.  The isotropic assumption is a
    first-order approximation; anisotropy would split the mode degeneracies
    of the equivalent sphere but would not change the order of magnitude of
    the eigenfrequencies or the coupling ratio.  Recent abdominal-wall
    mechanics reviews and elastography surveys place the canonical
    $E = \SI{0.1}{\mega\pascal}$ and $\eta = 0.25$ values as representative of
    contemporary measurement ranges, so the remaining uncertainty is
    quantitative rather than order-of-magnitude \cite{Deeken2017,Sack2023}.

  \item \textbf{Free-shell boundary conditions.}  The abdomen is partially
    constrained by the spine posteriorly, the pelvis inferiorly, and the
    rib cage superiorly.  The Rayleigh--Ritz study
    (\S\ref{sec:bc_effects}) is best understood as an \emph{exploratory
    sensitivity study}: it shows the range of possible frequencies under
    progressively constrained assumptions, from a free sphere
    ($f_2 = \SI{4.13}{\hertz}$) to 25\% surface clamped
    ($f_2 = \SI{3.66}{\hertz}$, $-11.4\%$).  Within the limited resolution
    of the present Ritz basis, these shifts suggest that fluid-dominated
    inertia may blunt the frequency impact of partial clamping, but that
    interpretation should be treated as illustrative rather than definitive
    until a formal convergence study is performed.  The entire bracketed range
    lies well within the material-property uncertainty band.  Boundary
    conditions do, however, affect the mechanical pathway through the modal
    participation factor $\Gamma_2$ (\S\ref{sec:coupling_ratio}), so the
    transfer-function separation remains large but is not literally unchanged.

  \item \textbf{Thin-shell approximation.}  The wall-thickness-to-radius
    ratio ($h/R \approx 0.06$) lies near the conventional upper bound of
    thin-shell theory ($h/R < 0.05$).  First-order shear-deformation
    (Mindlin--Reissner) corrections would reduce the $n = 2$ frequency by
    approximately \SI{4}{\percent} \cite{Soedel2004}, a negligible effect
    compared with the material-property uncertainty.  This limitation is
    unlikely to affect any conclusions.

  \item \textbf{Equivalent-sphere geometry.}  The oblate spheroidal
    abdomen is mapped to an equivalent sphere via $R_{\mathrm{eq}} =
    (a^2 c)^{1/3}$.  This preserves volume and leading-order curvature
    but suppresses the equatorial--polar mode splitting that occurs in
    an oblate spheroid \cite{Leissa1973}.  For a typical aspect ratio
    $c/a = 0.67$, the splitting is estimated at ${\sim}\SI{33}{\percent}$,
    which could shift individual mode frequencies but would not change the
    coupling ratio or the qualitative separation between breathing and
    flexural mode families.

    More fundamentally, the human abdominal cavity is a \emph{triaxial}
    ellipsoid (lateral ${\sim}\SI{14}{\centi\metre}$, craniocaudal
    ${\sim}\SI{16}{\centi\metre}$, anteroposterior
    ${\sim}\SI{10}{\centi\metre}$), not a spheroid of either kind.  The
    oblate approximation captures the principal feature---anterior--posterior
    flattening---but any spheroid is a simplification.  Crucially, however,
    the equivalent-sphere flexural frequencies depend on $R_{\mathrm{eq}}$
    alone: under the equivalent-sphere formula, at fixed $R_{\mathrm{eq}}$,
    varying $c/a$ from 0.3 (extremely oblate) to 3.0 (extremely prolate)
    changes $f_2$ by less than \SI{1}{\percent}.  The Ritz oblate-spheroid
    correction introduces additional aspect-ratio dependence (see
    Table~\ref{tab:oblate_correction}), but the geometry enters the coupling
    ratio primarily through $R_{\mathrm{eq}}$, so the qualitative results
    are insensitive to the choice of spheroidal idealisation.

  \item \textbf{Inviscid fluid assumption.}  The enclosed fluid is modelled
    as inviscid.  A viscous correction using the Stokes boundary-layer
    thickness shows that even at the worst-case mucus viscosity
    ($\mu = \SI{0.1}{\pascal\second}$), the quality factor decreases by only
    $\Delta Q = -1.85\%$.  The crossover viscosity at which viscous damping
    would equal structural damping is \SI{281}{\pascal\second}, far above
    physiologically realistic abdominal fluid viscosities.  The inviscid assumption is therefore
    well justified for any physiologically realistic abdominal fluid.

  \item \textbf{No internal organ structure.}  Visceral organs are
    homogenised into a single fluid.  This estimate follows from a
    Hashin--Shtrikman upper-bound effective-medium calculation for a fluid
    matrix containing soft solid inclusions \cite{Hashin1963}: taking
    $K_\mathrm{tissue} \approx K_\mathrm{water}$,
    $\mu_\mathrm{tissue} \approx \SI{1}{\kilo\pascal}$, and a volume fraction
    $\phi = 0.25$ (representative of liver, spleen, and kidneys combined)
    predicts a frequency shift of only \SI{0.4}{\percent} for the $n=2$ mode.
    Because the organ tissue bulk modulus is close to that of water, the
    effective compressibility barely changes, and at sub-percolation volume
    fractions the isolated inclusions cannot transmit a system-spanning shear
    network.  The homogeneous fluid model is therefore adequate for the
    present purposes.

    \item \textbf{Speculative local-coupling hypotheses.}  The brief discussion
    in \S\ref{sec:disc:alternatives} is intentionally schematic.  Any
    realistic assessment of gas pockets or orifice coupling would require a
    coupled gas--tissue model with anatomical confinement, wall loading, and
    transient boundary conditions, all of which are omitted here.

  \item \textbf{Mechanotransduction benchmarks.}  The
    \SIrange{0.5}{2.0}{\micro\metre} PIEZO activation range
    \cite{coste2010piezo1,Lewis2015,Xiao2024} derives from cell-culture and
    membrane-tension experiments and may not directly translate to
    tissue-scale effects, where strain concentrations around cells are
    modulated by the extracellular matrix.  We therefore treat these values as
    cellular mechanosensory reference scales rather than experimentally confirmed tissue-scale
    thresholds.  In gastrointestinal tissue the
    relevant transduction pathway is also likely to involve PIEZO2-expressing
    afferents rather than PIEZO1 alone \cite{ServinVences2023}.  This
    limitation introduces uncertainty in the absolute benchmark but does not
    affect the \emph{relative} comparison between airborne and mechanical
    coupling.

\end{enumerate}

\noindent
In summary, the identified limitations affect the absolute values of
predicted frequencies and displacements---in several cases substantially---but
none alter the central finding that airborne acoustic coupling to abdominal
flexural modes is $10^4$--$10^5$ times weaker than mechanical coupling under
occupationally relevant conditions.

\subsection{Proposed Experimental Validation}
\label{sec:disc:validation}

A useful validation sequence begins with a tissue-mimicking phantom having
the canonical spheroidal geometry and stiffness scale of
\S\ref{sec:formulation}.  Modal testing of such a phantom would check the
predicted flexural frequencies and mode shapes directly, while paired
mechanical and airborne loading would test the central claim of this paper:
that the coupling disparity between whole-body vibration and airborne
infrasound is orders of magnitude rather than marginal.  A natural in~vivo
extension would then combine subject-specific elastography with measured
abdominal transmissibility, allowing the same framework to be assessed
against human variability without changing the underlying model structure.
